\documentclass[aspectratio=169]{beamer}
\usepackage[T1,T2A]{fontenc}
\usepackage[utf8]{inputenc}
\usepackage[main=russian,english]{babel}
\usepackage{cmap,booktabs,amsmath,amssymb}
\usecolortheme{beaver}
\usefonttheme[onlymath]{serif}
\setbeamertemplate{navigation symbols}{}
\setbeamertemplate{footline}[page number]
\setbeamertemplate{blocks}[rounded=true,shadow=true]
\setbeamersize{text margin left=0.5cm,text margin right=0.5cm}
\setbeamerfont{author}{size=\normalsize}
\setbeamerfont{institute}{size=\normalsize}
\title{ReMix: градиенты для конечных смесей}
\subtitle{Отчёт по проекту}
\author{Александр Уденеев, Алексей Пономарев,\\ Лина Горбунова, Максим Иванов}
\institute{ИАД, МФТИ}
\date{}
\begin{document}
\begin{frame}
\thispagestyle{empty}
\maketitle
\end{frame}

\begin{frame}{Как обучать веса смеси?}
\[
 q_\theta(z)=\sum_{k=1}^K \pi_k q_{\phi_k}(z),
 \qquad \pi=\operatorname{softmax}(a).
\]
\begin{block}{Задача}
Подобрать веса $\pi_k$ и параметры компонент $\phi_k$, максимизируя
$\mathbb E_{q_\theta}[\log p(x,z)-\log q_\theta(z)]$.
\end{block}
Выбор номера компоненты дискретен. Обычный обратный проход через этот выбор не даёт градиента по весам.
\medskip

\textbf{ReMix} строит стохастическую оценку градиента и передаёт её оптимизатору PyTorch.
\end{frame}

\begin{frame}{Четыре способа оценить один градиент}
\small
\renewcommand{\arraystretch}{1.4}
\begin{tabular}{@{}p{.21\linewidth}p{.41\linewidth}p{.31\linewidth}@{}}
\toprule
Оценщик & Откуда берётся градиент & Что требуется \\
\midrule
\texttt{ScoreFunction} & Взвешивает $\nabla_\theta\log q_\theta(z)$ выигрышем точки & Выборка и дифференцируемая по параметрам лог-плотность \\
\texttt{Stratified} & Усредняет внутри каждой компоненты, затем по весам смеси & \texttt{rsample()} у компонент \\
\texttt{Implicit} & Сохраняет квантиль при изменении параметров смеси & Скалярные непрерывные компоненты с \texttt{rsample}, \texttt{cdf} \\
\texttt{PostStratified} & Усредняет вклады возможных меток при известной точке & Как у \texttt{Implicit}, плюс \texttt{icdf} \\
\bottomrule
\end{tabular}
\end{frame}

\begin{frame}{Из чего состоит один шаг обучения}
\begin{enumerate}
\item \texttt{Mixture(logits, components)} задаёт смесь распределений PyTorch и вычисляет $\log q_\theta(z)$.
\item \texttt{ELBOObjective(log\_joint)} задаёт вклад одной точки в цель.
\item \texttt{estimator.estimate(q, objective)} берёт выборку и возвращает \texttt{EstimateResult}.
\end{enumerate}
\begin{columns}[T,onlytextwidth]
\column{.46\textwidth}
\begin{block}{\texttt{result.value}}
Оценка цели для журнала обучения. Тензор отделён от графа.
\end{block}
\column{.51\textwidth}
\begin{block}{\texttt{result.surrogate\_loss}}
\texttt{backward()} вычисляет градиент для шага минимизации в PyTorch.
\end{block}
\end{columns}
\end{frame}

\begin{frame}{Зачем выделять Objective в отдельный класс?}
\texttt{ELBOObjective.evaluate(q, z)} возвращает по одному значению
$\log p(x,z)-\log q_\theta(z)$ на точку и проверяет совпадение форм.
При \texttt{stl=True} параметры $q$ внутри $\log q$ отсоединяются от графа, путь через $z$ сохраняется.
\begin{block}{Можно ли заменить ELBO?}
Для цели $\mathbb E_q[f(q,z)]$ можно передать свой объект с методом \texttt{evaluate(q, z)}. Например, для $\beta$-VAE при фиксированном $x$:
\[
 f=\log p(x\mid z)+\beta\,[\log p(z)-\log q(z\mid x)].
\]
\end{block}
Формулу и правила дифференцирования можно менять в одном месте для всех четырёх оценщиков. Для другой дивергенции нужно проверить вид матожидания и условия выбранного метода.
\end{frame}

\begin{frame}{Нужен ли смеси свой \texttt{rsample()}?}
\begin{block}{Что есть в версии 1.0.0}
\texttt{Mixture.sample()} выбирает компоненту и возвращает точку без графа. У самой смеси метода \texttt{rsample()} нет.
\end{block}
\begin{itemize}
\item В \texttt{Stratified} градиент проходит через \texttt{rsample()} каждой компоненты и через веса их средних.
\item В \texttt{Implicit} производная функции распределения смеси входит прямо в оценку градиента цели.
\end{itemize}
Общий \texttt{Mixture.rsample()} потребовал бы указать, для каких смесей определён такой градиент.
\end{frame}

\begin{frame}{Результат проекта}
\begin{block}{\texttt{remix-vi 1.0.0}}
Опубликован пакет для PyTorch. В одном цикле обучения можно заменить оценщик, сохранив смесь, целевую функцию и оптимизатор.
\end{block}
Выбор метода зависит от доступных операций над компонентами и размерности события.
\medskip

\href{https://pypi.org/project/remix-vi/1.0.0/}{\beamergotobutton{Страница релиза на PyPI}}
\end{frame}
\end{document}
